% !PRINTED 112
\weberpage{112}
In the field $R(i)$ we have (Bd. I, \S.\,173) the units $\pm 1, \pm i$, further,
as primes, the associated factors $1\pm i$ of 2, the real primes
$q$ of the form $4N+3$, and the two complex factors $\pi, \pi'$ of the real
primes $p$ of the form $4N+1$.

In Band I, \S.\,172 we found the decomposition

\[
p = \psi_1(i)\,\psi_1(-i)
\]

which permits us to set

\weq{10}{\pi=\psi_1(i), \quad \pi'=\psi_1(-i)}

and, if $\eta$ is an $N$-term period of $p$th
roots of unity,

\[
(i,\eta)^4 = p\,\psi_1(i)^2, \quad (-i,\eta)^4 = p\,\psi_1(-i)^2,
\]

or

\weq{11}{(i,\eta)^4=\pi^3\pi', \quad (-i,\eta)^4=\pi\pi'^3, \quad (i,\eta)(-i,\eta)=p.}

If we now decompose the numbers $a+bi$, $a-bi$ into their prime factors
in $R(i)$, as with the cubic equations, then, employing a corresponding
notation, we obtain

\weq{12}{\begin{aligned}
a+bi &= i^\lambda (1+i)^n q_1^{s_1} q_2^{s_2} \ldots \pi_1^{t_1} \pi_1'^{t_1'} \pi_2^{t_2} \pi_2'^{t_2'} \ldots \\
a-bi &= i^{-\lambda}(1-i)^n q_1^{s_1} q_2^{s_2} \ldots \pi_1^{t_1'} \pi_1'^{t_1} \pi_2^{t_2'} \pi_2'^{t_2} \ldots
\end{aligned}}

But the product of these two numbers must be the fourth power of a rational
number [by (8) and (9)], and from this it follows:

\weq{13}{\begin{aligned}
n &\equiv 0 & &\text{(mod 4)} \\
s_1 &\equiv 0, \quad s_2 \equiv 0 \ldots & &\text{(mod 2)} \\
t_1+t_1' &\equiv 0, \quad t_2+t_2' \equiv 0 \ldots & &\text{(mod 4)} \\
t_1-t_1' &\equiv 0, \quad t_2-t_2' \equiv 0 \ldots & &\text{(mod 2).}
\end{aligned}}

But now it follows from (11):

\weq{14}{\pi^t\pi'^{t'} = (\pi^3\pi')^{\frac{t-t'}{2}}(\pi\pi')^{\frac{-t+3t'}{2}} = (i,\eta)^{2(t-t')}\,p^{\frac{-t+3t'}{2}},}

and further

\weq{15}{(1+i)^n = (-1)^{\frac{n}{4}}\,2^{\frac{n}{2}},}

and accordingly, having regard to (13), we can therefore, if we denote by $H_1, H_2$
two cyclotomic numbers, namely the product of all the numbers occurring in the
decomposition of $a+bi$

\[
(i,\eta)^{\frac{t-t'}{2}},
\]
%CONT>>
